\documentclass[11pt]{article}
\usepackage[utf8]{inputenc}\usepackage[T1]{fontenc}
\usepackage[spanish,es-nodecimaldot,es-noshorthands]{babel}
\usepackage[margin=1in]{geometry}
\usepackage{amsmath,amssymb,booktabs}
\newcommand{\E}{\mathbb E}
\usepackage[colorlinks=true,allcolors=blue!50!black]{hyperref}
\title{Guía completa para escribir las derivaciones a mano}
\author{Material de apoyo para Fabián Matías Galarza Chumbe}
\date{10 de noviembre de 2026}
\begin{document}\maketitle
\noindent\textbf{Estado: guía tipográfica, no manuscrito del autor.}
Cada paso siguiente debe escribirse y revisarse personalmente. Al digitalizar,
incluir nombre, fecha y número de página. Usar tinta legible, sin sustituir
una fuente tipográfica por supuesta escritura humana. La revisión de
comprensión todavía no ha ocurrido en esta sesión.

\section{01: objetos, unidades y dominio}
Escribir primero qué se mantiene fijo: $\kappa$, parámetros de tecnología,
precio de pares y horizonte del evento. Definir $D$ como tasa anual decimal,
$q^0,q^*$ como probabilidades, $e,\kappa$ como índices y pagos por sol de crédito.
\begin{enumerate}
\item $q^0=\bar q+aD$, $a>0$, $0<\bar q<1$.
\item $0\le q^0\le1$ equivale a $-\bar q\le aD\le1-\bar q$ por restar $\bar q$.
\item Dividir por $a>0$ conserva desigualdades:
$-\bar q/a\le D\le(1-\bar q)/a$.
\item $R=\bar R+D>0$ equivale a $D>-\bar R$.
\item Elegir el dominio $\mathcal D$ dentro de la intersección anterior.
No extrapolar la tecnología afín fuera de ese dominio.
\item $q=q^0-ge$, con $g>0$. Como $e\ge0$, $q\le q^0\le1$.
\item Para $q\ge0$, despejar $q^0-ge\ge0$, luego $ge\le q^0$ y
$e\le q^0/g$ dividiendo por $g>0$.
\item Combinar con capacidad $e\le\bar e\ge0$:
$0\le e\le u=\min\{\bar e,q^0/g\}$.
\end{enumerate}
Anotar al margen que una probabilidad de default de por vida y un saldo con
atraso a fecha $h$ no son intercambiables sin el Supuesto 5.

\section{02: señal de precios y cancelación exacta}
Escribir $n$ pares distintos de $j$, $n\ge3$ para el diseño propuesto:
\begin{align*}
D&=(i_j-b)-\frac1n\sum_{k\ne j}(i_k-b)\\
 &=i_j-b-\left(\frac1n\sum_k i_k-\frac1n\sum_k b\right)\\
 &=i_j-b-\frac1n\sum_k i_k+b\\
 &=i_j-\frac1n\sum_k i_k.
\end{align*}
Reglas: distributividad, linealidad de la suma, $\sum_k b=nb$ porque la base
es idéntica para todos. Para mediana, escribir que trasladar cada observación
por la misma constante traslada su mediana en esa constante. No es una
igualdad entre promedio y mediana.

Para cargas heterogéneas, sustituir $i_j=b+\lambda_j z+v_j$ y el promedio
$b+\bar\lambda z+\bar v$. Restar término a término:
$D=(\lambda_j-\bar\lambda)z+(v_j-\bar v)$. Esto demuestra por qué el ciclo
no se cancela si las cargas difieren.

Ejemplo hipotético, porcentajes: $(8+9+10+25)/4=52/4=13$; mediana
$(9+10)/2=9.5$; $12-13=-1$ pp y $12-9.5=2.5$ pp.

\section{03: derivación de la FOC y la concavidad}
Suprimir $\kappa$ únicamente después de definirlo. El objetivo es
\[
\pi(e)=(1-q^0+ge)R-F-\frac c2e^2.
\]
Mantener $D$, y por tanto $q^0$ y $R$, fijos al derivar respecto de $e$.
\begin{enumerate}
\item Distribuir el pago: $\pi=(1-q^0)R+gRe-F-(c/2)e^2$.
\item La derivada de $(1-q^0)R$ es cero porque no depende de $e$.
\item La derivada de $gRe$ es $gR$ por regla de potencia y coeficiente constante.
\item La derivada de $-F$ es cero.
\item La derivada de $-(c/2)e^2$ es $-(c/2)2e=-ce$.
\item Sumar: $\pi_e=gR-ce$.
\item Derivar otra vez: $\pi_{ee}=-c<0$ porque $c>0$.
\item Una solución interior satisface $gR-ce=0$.
\item Sumar $ce$ a ambos lados: $gR=ce$.
\item Dividir por $c>0$: $z=e=gR/c>0$ porque $g,R>0$.
\end{enumerate}
La FOC identifica un candidato no restringido. Todavía no se ha comprobado
que $z\le u$, por lo que no es aún la solución global.

\section{04: cuadrado completo, existencia y unicidad (P1)}
Desarrollar explícitamente el cuadrado:
\begin{align*}
-\frac c2(e-z)^2
 &=-\frac c2(e^2-2ez+z^2)\\
 &=-\frac c2e^2+cez-\frac c2z^2\\
 &=-\frac c2e^2+gRe-\frac{g^2R^2}{2c},
\end{align*}
porque $z=gR/c$, $cz=gR$ y $cz^2/2=g^2R^2/(2c)$. Sumar
$(1-q^0)R-F+g^2R^2/(2c)$ a la expresión anterior recupera exactamente $\pi$.

Condiciones de existencia: $u\ge0$, intervalo $[0,u]$ compacto no vacío y
objetivo continuo. Por el teorema del valor extremo hay máximo. Si $u=0$,
solo es factible $e=0$. Si $u>0$, $-c<0$ da concavidad estricta, por tanto
a lo sumo un máximo en el intervalo convexo.

Ahora minimizar $(e-z)^2$, con $z>0$. Si $0<z<u$, su valor cero ocurre solo
en $e=z$. Si $z\ge u$, para cualquier $e<u$ se cumple
$z-e>z-u\ge0$; elevar al cuadrado conserva el orden estricto, así que
$(e-z)^2>(u-z)^2$. El único máximo es $u$. Si $z=u$, sigue siendo único.
En consecuencia
\[
e^*=\min\{z,u\}=\min\{gR/c,\bar e,q^0/g\}.
\]
No hay esquina inferior con capacidad positiva y $q^0>0$, porque $\pi_e(0)=gR>0$.
Si $q^0=0$ o $\bar e=0$, el intervalo es singleton y sí se elige cero.

\section{05: sustitución, pisos y pendientes (P1)}
Escribir $v=\min\{\bar e,gR/c\}\ge0$.
\begin{enumerate}
\item $e^*=\min\{v,q^0/g\}$.
\item Si $v\le q^0/g$, entonces $e^*=v$ y $q^*=q^0-gv\ge0$.
\item Si $v>q^0/g$, entonces $e^*=q^0/g$ y $q^*=q^0-g(q^0/g)=0$.
\item Ambos casos equivalen a $q^*=\max\{0,q^0-gv\}$.
\end{enumerate}
Derivar cada régimen únicamente en un intervalo donde las desigualdades son
estrictas y se mantienen $\kappa,a,g,c,\bar e$ fijos:
\begin{align*}
\text{Interior: }q^*&=\bar q+aD-(g^2/c)(\bar R+D),\\
\frac{dq^*}{dD}&=0+a-(g^2/c)(0+1)=a-g^2/c,\\
\frac{de^*}{dD}&=\frac{g}{c}\frac{d(\bar R+D)}{dD}=g/c>0;\\
\text{Capacidad: }q^*&=\bar q+aD-g\bar e,\quad \frac{dq^*}{dD}=a;\\
\text{Piso: }q^*&=0,\quad \frac{dq^*}{dD}=0.
\end{align*}
En un empate escribir las fórmulas a izquierda y derecha. La solución
continúa porque mínimo y máximo conservan continuidad; la derivada existe
solo si las pendientes laterales coinciden. No confundir continuidad con
diferenciabilidad.

Ejemplo de capacidad: $\bar q=.35,a=.04,g=.4,c=1,\bar e=.44,\bar R=1.10$.
Igualar $g(\bar R+D)/c=\bar e$ da $D=c\bar e/g-\bar R=1.10-1.10=0$.
A la izquierda $a-g^2/c=.04-.16=-.12$; a la derecha $a=.04$.
Comprobar en $D=0$: $q^0/g=.35/.4=.875>.44$, por lo que la cota de
probabilidad no causa ese quiebre. Riesgo en el punto: $.35-.4(.44)=.174$.

\section{06: umbral y condiciones de su existencia (P2)}
\begin{enumerate}
\item Definir $H(\kappa)=g(\kappa)^2/c(\kappa)>0$ y $\beta=a-H$.
\item Continuidad de $a,g,c$ y $c>0$ implica continuidad de $H$ por producto
y cociente, y de $\beta$ por resta.
\item Comprobar $\beta(0)>0$ y $\beta(1)<0$ antes de aplicar valor intermedio.
\item El teorema da un $\kappa^*\in(0,1)$ con $\beta(\kappa^*)=0$.
\item Si $x<y$ fueran dos ceros, monotonía estricta da $\beta(x)>\beta(y)$,
pero ambos son cero: contradicción. Hay un único cero.
\item Para $\kappa<\kappa^*$, $\beta(\kappa)>\beta(\kappa^*)=0$.
Para $\kappa>\kappa^*$, $\beta(\kappa)<0$.
\item La identificación de $\beta$ con $\partial_Dq^*$ requiere verificar
interioridad en todo $\kappa$ al precio $D_0$ comparado.
\end{enumerate}
No se requiere diferenciar en $\kappa$ para probar esta versión. Monotonía
débil no garantiza unicidad. Continuidad sola tampoco garantiza un cero si
no hay cruce ni un signo cero en extremos.

\section{07: logaritmos, calibración y alternativa}
Como $H,g,c>0$, los logaritmos existen. Por reglas de logaritmos,
$\log H=2\log g-\log c$. Por regla de la cadena:
$H'/H=2g'/g-c'/c$. Si $g'\le0$ y $c'\ge0$, el lado derecho es no positivo,
por tanto $H'\le0$ al multiplicar por $H>0$.

Para los parámetros ilustrativos:
\begin{align*}
H&=\frac{(\sqrt{.05}e^{.8\kappa})^2}{e^{.4\kappa}}
=.05e^{1.6\kappa-.4\kappa}=.05e^{1.2\kappa},\\
a=H&\iff .18e^{-.8\kappa}=.05e^{1.2\kappa}\\
&\iff 3.6=e^{2\kappa}\iff \log3.6=2\kappa
\iff \kappa^*=\log(3.6)/2.
\end{align*}
Justificar cada división por un número positivo, el producto de
exponenciales y que $\log$ es inversa de la exponencial en los positivos.
El valor es aproximadamente $.640467$. Sustituir extremos en $a-H$:
$.18-.05=.13>0$ y $.18e^{-.8}-.05e^{1.2}\approx-.085127<0$.

En la alternativa $g=\sqrt{.05}e^{-.4\kappa}$,
$H=.05e^{-1.2\kappa}$ y $a/H=3.6e^{.4\kappa}>1$ para $\kappa\in[0,1]$.
Como $H>0$, $a-H=H(a/H-1)>0$. No hay umbral. Escribir por qué este
contraejemplo satisface probabilidades y esfuerzo admisibles en la simulación.

\section{08: contabilidad de cohortes y esperanza (P3)}
Definir $B_0,m_0,B,q^*$ antes de usar los símbolos. Bajo saldo determinista o
población continua, atraso esperado antiguo $B_0m_0$, nuevo $Bq^*$, total
$A=B_0m_0+Bq^*$. Saldo total $T=B_0+B>0$.
\begin{align*}
M&=A/T=\frac{B_0}{B_0+B}m_0+\frac{B}{B_0+B}q^*,\\
w&=B/(B_0+B),\quad
1-w=(B_0+B-B)/(B_0+B)=B_0/(B_0+B),\\
M&=(1-w)m_0+wq^*.
\end{align*}
Si el denominador fuera aleatorio, $\E[A/T]$ no se reemplaza en general por
$\E[A]/\E[T]$. Escribir ese límite y el evento de atraso al horizonte. No
confundir tasa de incidencia de impago con fracción de saldo en atraso.

\section{09: derivada, conservación e inversión del signo (P3)}
Con $m_0$ fijo y funciones diferenciables en el régimen:
\begin{align*}
M'&=(1-w)'m_0+(1-w)m_0'+w'q^*+wq^{*\prime}\\
&=-w'm_0+0+w'q^*+w\beta_D\\
&=w\beta_D+w'(q^*-m_0).
\end{align*}
Reglas: producto, derivada de constante, agrupación distributiva.
Si $w'=0$, $M'=w\beta_D$ y $w>0$ preserva signos.
Si $\beta_D>0$ y $|C|<w\beta_D$ con $C=w'(q^*-m_0)$, entonces
$M'=w\beta_D+C\ge w\beta_D-|C|>0$.
Si $\beta_D<0$, $M'\le w\beta_D+|C|<0$.
La condición estricta no cubre $\beta_D=0$ ni empates de magnitudes.

Derivar el peso sin omitir el cociente:
\begin{align*}
w'&=\frac{B'(B_0+B)-B(0+B')}{(B_0+B)^2}
=\frac{B'B_0+B'B-BB'}{(B_0+B)^2}
=\frac{B'B_0}{(B_0+B)^2}.
\end{align*}
Para $B=B_1e^{-\varepsilon D}$, regla de cadena:
$B'=B_1e^{-\varepsilon D}(-\varepsilon)=-\varepsilon B$.
Luego
\[
w'=-\varepsilon\frac{B}{B_0+B}\frac{B_0}{B_0+B}
=-\varepsilon w(1-w).
\]
Sustituir en $M'$ y factorizar $w>0$.

Si $\beta_D>0$ y $q^*>m_0$, escribir
$M'<0\iff\beta_D-\varepsilon(1-w)(q^*-m_0)<0
\iff\varepsilon(1-w)(q^*-m_0)>\beta_D$.
No se cambia el signo al dividir por $w$, que es positivo.

\section{10: contraejemplo exacto y unidades}
En $D=0$ y $\kappa=0$, $q^0=.28$, $H=.05$, $R=1.10$.
\begin{align*}
q^*&=.28-.05(1.10)=.28-.055=.225=9/40,\\
\beta_D&=.18-.05=.13=13/100,\\
w&=1/(4+1)=.2=1/5,\quad m_0=.10=1/10,\\
q^*-m_0&=.125=1/8,\\
\varepsilon(1-w)(q^*-m_0)&=3(4/5)(1/8)=3/10=.30,\\
M'&=(1/5)(13/100-30/100)=-17/500=-.034.
\end{align*}
Comprobar $0<q^*,m_0,w<1$; $B_0=4>0$, $B_1=1>0$ y $\varepsilon=3\ge0$.
En las primitivas originales $g=\sqrt{.05},c=1$,
$e^*=\sqrt{.05}(1.10)\approx.245967<.50$ y
$q^0/g\approx1.2522>e^*$: interioridad real, no supuesta por la conclusión.

Una unidad de $D$ es cien puntos porcentuales. Para $\Delta D=.01$,
el cambio local aproximado de $q$ es $.13(.01)=.0013$, es decir $.13$ pp;
el de $M$ es $-.034(.01)=-.00034$, es decir $-.034$ pp. Usar el símbolo
de aproximación para cambios finitos. Lean comprueba la aritmética racional
de estas expresiones, no sus derivadas ni los supuestos económicos.

\section{11: comprobaciones y qué escribir junto a cada prueba}
Al final de cada hoja anotar: resultado que se probó, supuestos usados, regla
matemática aplicada, caso de esquina y ejemplo que falsaría una versión más
fuerte. Para simulaciones, describir sin presentarlas como demostraciones:
\begin{enumerate}
\item Optimización independiente: evaluar el objetivo y ambos extremos.
\item Diferencias centrales: $[q(D+\delta)-q(D-\delta)]/(2\delta)$,
con $\delta=10^{-6}$, verificando mismo régimen en los tres puntos.
\item Error de esfuerzo máximo comparado con tolerancia $2\times10^{-6}$.
\item No extrapolar una malla de parámetros al espacio continuo completo.
\end{enumerate}
La fórmula de error numérico es $|e_{\text{analítico}}-e_{\text{optimizador}}|$.
La prueba de existencia, unicidad y signo usa argumentos analíticos del
artículo, no la ausencia de fallos en una malla.

\section{12: control personal y entrega}
Digitalizar las hojas en un PDF con páginas numeradas o en archivos por tema.
Por ejemplo: \path{01-dominio.pdf}, \path{02-precio.pdf} y
\path{10-contraejemplo.pdf}.

Incluir los casos $u=0$, empate de cotas, derivadas laterales, conservación
de signo y alternativa sin umbral. Si cambia la numeración, actualizar las
referencias del artículo y las diapositivas.

Antes de afirmar comprensión, responder sin leer: ¿qué se mantiene fijo?,
¿por qué el máximo es único?, ¿por qué evaluar reduce riesgo sin cambiar
volumen en este modelo?, ¿cuándo falla el umbral?, ¿por qué un ratio de saldo
puede dar el signo contrario?, ¿qué no ha comprobado Lean?, ¿qué tendría que
cambiar para una interpretación causal? Si una respuesta no sale con claridad,
volver a la hoja correspondiente y corregirla.
\end{document}
